\section{Выполнение программы}

\subsection{Подготовка к работе}

Перед запуском программы оператор должен:
\begin{enumerate}
    \item убедиться в наличии исполняемого файла программы;
    \item определить каталог или URI вывода результатов;
    \item выбрать формат хранения выходных данных;
    \item определить масштаб генерации и перечень таблиц;
    \item проверить наличие прав на запись в целевой каталог или хранилище.
\end{enumerate}

\subsection{Загрузка программы}

Программа не требует отдельной процедуры загрузки в память оператором и
запускается как обычное консольное приложение из командной строки.

Общий вид вызова программы:

\begin{verbatim}
schengen_main [параметры]
\end{verbatim}

\subsection{Запуск программы}

Для получения справки оператор должен использовать команду:

\begin{verbatim}
schengen_main --help
\end{verbatim}

Для вывода списка доступных таблиц используется команда:

\begin{verbatim}
schengen_main --list-tables
\end{verbatim}

Для вывода списка доступных форматов используется команда:

\begin{verbatim}
schengen_main --list-formats
\end{verbatim}

Пример типового запуска программы:

\begin{verbatim}
schengen_main \
  --scale-factor 1 \
  --output-format parquet \
  --output-path ./out \
  --table customer orders lineitem
\end{verbatim}

\subsection{Команды оператора}

Основными параметрами командной строки являются:
\begin{enumerate}
    \item \texttt{--help} -- вывод справки;
    \item \texttt{--list-tables} -- вывод списка доступных таблиц;
    \item \texttt{--list-formats} -- вывод списка доступных форматов;
    \item \texttt{--scale-factor} -- задание масштаба генерации;
    \item \texttt{--output-format} -- задание формата выходных данных;
    \item \texttt{--output-path} -- задание каталога или URI вывода;
    \item \texttt{--table} -- указание одной или нескольких таблиц для генерации;
    \item \texttt{--batch-rows} -- задание числа строк в батче;
    \item \texttt{--write-strategy} -- выбор стратегии записи;
    \item \texttt{--write-workers} -- задание числа рабочих потоков обработки;
    \item \texttt{--write-queue-capacity} -- задание емкости очереди записи;
    \item \texttt{--verbosity} -- выбор режима вывода хода выполнения.
\end{enumerate}

Для формата \texttt{parquet} могут использоваться дополнительные параметры:
\begin{enumerate}
    \item \texttt{--parquet-use-threads};
    \item \texttt{--parquet-row-group-bytes};
    \item \texttt{--parquet-max-row-group-rows};
    \item \texttt{--parquet-output-buffer-bytes};
    \item \texttt{--parquet-compression}.
\end{enumerate}

Для формата \texttt{orc} могут использоваться дополнительные параметры:
\begin{enumerate}
    \item \texttt{--orc-max-partition-rows};
    \item \texttt{--orc-stripe-bytes};
    \item \texttt{--orc-max-stripe-rows};
    \item \texttt{--orc-output-buffer-bytes};
    \item \texttt{--orc-compression}.
\end{enumerate}

\subsection{Управление выполнением программы}

Если параметр \texttt{--verbosity verbose} задан, программа выводит сведения о ходе
выполнения по таблицам. Если параметр \texttt{--verbosity quiet} не изменен, программа
работает в минимальном режиме вывода и сообщает оператору преимущественно об
ошибках.

Стратегия записи задается параметром \texttt{--write-strategy}. Допустимыми значениями
являются:
\begin{enumerate}
    \item \texttt{auto};
    \item \texttt{parallel-files};
    \item \texttt{single-file-ordered}.
\end{enumerate}

Число потоков обработки задается параметром \texttt{--write-workers}. Значение \texttt{0}
означает автоматический выбор.

\subsection{Завершение работы программы}

После завершения генерации программа закрывает используемые потоки обработки,
завершает запись выходных данных и возвращает код завершения в операционную
систему.

При успешном завершении программа возвращает код \texttt{0}. При ошибке параметров
запуска, неподдерживаемом формате, недопустимой стратегии записи либо сбое
выполнения программа возвращает ненулевой код завершения.

\subsection{Примеры выполнения}

Пример генерации всех таблиц в формат \texttt{parquet}:

\begin{verbatim}
schengen_main --scale-factor 1 --output-format parquet --output-path ./out
\end{verbatim}

Пример генерации выбранных таблиц в формате \texttt{orc}:

\begin{verbatim}
schengen_main \
  --scale-factor 1 \
  --output-format orc \
  --output-path ./out_orc \
  --table customer orders lineitem
\end{verbatim}

Пример запуска с подробным выводом хода выполнения:

\begin{verbatim}
schengen_main \
  --scale-factor 1 \
  --output-format parquet \
  --output-path ./out \
  --verbosity verbose
\end{verbatim}
